% Part: first-order-logic
% Chapter: models-theories
% Section: introduction

\documentclass[../../../include/open-logic-section]{subfiles}

\begin{document}

\olfileid[fr]{fol}{mat}{int}
\olsection{Introduction}

\begin{explain}
L'élaboration de la méthode axiomatique est un accomplissement
majeur de l'histoire des sciences, et revêt une importance
particulière dans l'histoire des mathématiques. Développer une
discipline par la méthode axiomatique exige de clarifier bien des
questions : quel est son objet ? Quels sont ses concepts les plus
fondamentaux ? Comment sont-ils liés ? Peut-on définir tous les
concepts de la discipline à partir de ces concepts fondamentaux ?
À quelles lois obéissent-ils et doivent-ils obéir ?

La méthode axiomatique et la logique se complètent naturellement.
La logique formelle fournit les outils nécessaires pour formuler des
théories axiomatiques, démontrer leurs théorèmes à partir des axiomes
suivant des règles précisément définies et étudier systématiquement
les propriétés de tous les systèmes qui satisfont ces axiomes.
\end{explain}

\begin{defn}
Un ensemble d'!!{sentence}s~$\Gamma$ est \emph{clos} si, chaque fois
que $\Gamma \Entails !A$, on a $!A \in \Gamma$. La \emph{clôture}
d'un ensemble d'!!{sentence}s~$\Gamma$ est
$\Setabs{!A}{\Gamma \Entails !A}$.

On dit que~$\Gamma$ est \emph{axiomatisé par} un ensemble
d'énoncés~$\Delta$ si $\Gamma$ est la clôture de~$\Delta$.
\end{defn}

\begin{explain}
Nous pouvons considérer une théorie axiomatique comme l'ensemble des
énoncés axiomatisé par son ensemble d'axiomes~$\Delta$. Autrement
dit, supposons donné un langage du premier ordre dont les symboles
non logiques désignent les notions primitives de la discipline que
nous voulons développer axiomatiquement, ainsi qu'un ensemble
d'!!{sentence}s qui expriment ses lois fondamentales. Nous pouvons
alors représenter la théorie par l'ensemble de tous les énoncés de
ce langage qui sont des conséquences des axiomes. Cela vaut aussi
bien pour des exemples simples, avec une seule notion primitive et
des axiomes simples, comme la théorie des ordres partiels, que pour
des théories complexes comme la mécanique newtonienne.

Voici les faits logiques qui rendent cette approche formelle de la
méthode axiomatique si importante. Supposons que $\Gamma$ soit un
système d'axiomes pour une théorie, c'est-à-dire un ensemble
d'énoncés.
\begin{enumerate}
\item Nous pouvons énoncer précisément à quelle condition un système
  d'axiomes caractérise une classe de !!{structure}s visée. Si une
  certaine classe de !!{structure}s nous intéresse, le système
  d'axiomes~$\Gamma$ la caractérise si et seulement si elle comprend
  exactement les !!{structure}s~$\Struct M$ telles que
  $\Sat{M}{\Gamma}$.
\item Cela peut échouer parce qu'il existe des $\Struct M$ telles
  que $\Sat{M}{\Gamma}$ mais que $\Struct M$ ne fait pas partie des
  !!{structure}s visées. Nous pouvons alors ajouter des axiomes qui
  ne sont pas vrais dans~$\Struct M$.
\item Si, au moins, $\Gamma$ est vrai dans toutes les !!{structure}s
  visées, alors tout énoncé~$!A$ tel que $\Gamma \Entails !A$ y est
  vrai également. Des outils logiques, tels que les méthodes de
  !!{derivation}, nous permettent donc d'établir qu'un énoncé est
  vrai dans toutes les !!{structure}s visées en montrant simplement
  qu'il découle des axiomes.
\item Parfois, nous ne partons d'aucune !!{structure} visée et
  commençons par les axiomes eux-mêmes : nous posons des notions
  primitives auxquelles nous voulons imposer certaines lois, codifiées
  par un système d'axiomes. Nous souhaitons vérifier d'emblée que
  les axiomes ne se contredisent pas. Dans le cas contraire, aucun
  concept ne peut obéir à ces lois et la théorie proposée est
  incohérente. Pour écarter cette possibilité, nous pouvons trouver
  un modèle de~$\Gamma$. S'il existe des modèles de notre théorie,
  nous pouvons aussi les étudier et en construire à l'aide de
  méthodes logiques.
\item L'indépendance des axiomes est également une question
  importante. L'un d'eux peut être une conséquence des autres et se
  révéler ainsi redondant. Pour démontrer qu'un axiome $!A$ de
  $\Gamma$ est redondant, nous pouvons établir
  $\Gamma \setminus \{!A\} \Entails !A$. Pour montrer qu'il ne
  l'est pas, nous pouvons aussi établir que
  $(\Gamma \setminus \{!A\}) \cup \{\lnot !A\}$ est satisfaisable.
  C'est ainsi, par exemple, qu'on a démontré l'indépendance du
  postulat des parallèles par rapport aux autres axiomes de la
  géométrie.
\item Une autre question importante concerne la définissabilité
  des concepts dans une théorie. Le choix du langage détermine de
  quoi se composent les modèles de la théorie. Toutefois, chaque
  aspect de celle-ci n'a pas nécessairement besoin d'y être
  représenté séparément. Par exemple, tout ordre~$\le$ détermine
  l'ordre strict correspondant~$<$ : dès que l'un est donné, nous
  pouvons définir l'autre. Un modèle d'une théorie comportant un
  tel ordre n'a donc pas besoin de contenir \emph{aussi} l'ordre
  strict correspondant. À quelles conditions, en général, une
  relation peut-elle être définie à partir d'autres ? Quand est-ce
  impossible, de sorte qu'il faut l'ajouter aux notions primitives
  du langage ?
\end{enumerate}
\end{explain}


\end{document}
